% sections/discussion.tex — numbers only through \res{}; every claim has a row in paper/CLAIMS.md.
% Draft for the author's review (plan of sections approved 2026-10-18, PLAN §6); revised after the review r1.
\section{Discussion}\label{sec:discussion}

\subsection{When does a QUBO pipeline pay?}\label{sec:disc-when}

The answer depends on the instance, on the sampler and on what a method starts from. With a cheap classical start ---
the continuous relaxation rounded to its $K$ largest weights and refined --- most of our factor and OR-Library
instances and a substantial share of the S\&P instances need no combinatorial search at all, and on the factor
instances the gaps that remain are a few basis points of certainty equivalent per year (Section~\ref{sec:res-a3}).
On the remaining instances, a pipeline whose sampler preserves the cardinality approached the optimum faster than the
open-source exact solvers, while the pipelines with the tuned penalty samplers did not (Section~\ref{sec:res-rq2}).
Against a local search on the true objective, the pipeline was faster on the typical factor instance but failed to
reach the target within the budget on a minority of them and on most of the S\&P instances; local search never
failed, and on average it was faster (the direction of H2b is established, its size relative to the band is not).
The pipeline therefore offers speed on many instances without reliability, and local search offers both on average.

A niche for a fast first portfolio is not supported either. Without the common start the pipeline holds its first
portfolio sooner than local search and the exact methods, which reverses the order of H2b (Table~\ref{tab:exact}), but
the common start itself is ready earlier on every instance, and the first portfolio of the pipeline is at least as
good as it on most, but not all, instances. For the CCMV problem, where the continuous relaxation is always available and cheap, a QUBO pipeline
did not deliver a portfolio sooner than the classical start.

The results of RQ1 suggest a mechanism, which we did not test at the sizes of RQ2. At $N \le 24$ the ceiling of the
equal-weight surrogate rises with the dispersion of volatilities, an oracle pool of the best subsets of the surrogate
recovers most of the loss, and the real samplers either miss the minimum of the surrogate (penalty annealing) or
find it but return few distinct subsets per call (tabu search and swap-SA), so that their pools recover less of the
loss than the oracle pool (Table~\ref{tab:samplers-rq1}). A pipeline that ranks subsets by the surrogate inherits its
ceiling unless its pool is diverse; local search ranks supports by the true objective at every step.

These findings bear on how QUBO approaches to portfolio selection are usually evaluated. Positive results in the
literature rest mostly on closed hybrid services and on weak or untuned classical
baselines~\cite{mugel2022dynamic,lang2022strategic,kim2025quantum}, while strong classical solvers, where they are
present, solve the portfolio problems considered~\cite{lozano2026where,stopfer2025quantum,chen2025benchmarking}. Our
benchmark differs from these studies in three respects: an anytime measure under equal budgets and one tuning
protocol for all methods~\cite{barr1995designing,bucher2024towards}, a local search on the true objective as a
classical baseline of the same kind as the pipeline, and a screening of trivial instances. On the calibration
instances the common start was already within the target on most instances, so without the screening the outcome
``equivalent'' could have been fixed by the composition of the cells rather than by the methods.

\subsection{What the exact baseline means for H2}\label{sec:disc-baseline}

On the problems of Bertsimas and Cory-Wright~\cite{bertsimas2022scalable}, our open-source exact representative was more
than an order of magnitude slower than the best published methods, while outer approximation with an open-source master
problem stayed within an order of magnitude on BCW-T3, and on the hardest instances of BCW-T4 SCIP was competitive
(Section~\ref{sec:res-m11}). The validation measures the time to a certificate, and H2 measures the speed of improving
an incumbent; the two need not agree, and the representative was chosen on the second. Its tuning space was small
--- the formulation and the symmetry handling, \res{e2.tune.exact.n100.nconfigs} and \res{e2.tune.exact.n200.nconfigs}
configurations searched completely, against \res{e2.tune.heur.nconfigs} random configurations of each heuristic.
Two post hoc runs probed this choice: with SCIP's aggressive primal heuristics, and with outer approximation as the
representative in every cell, the outcome of H2 was unchanged, and the aggressive heuristics did not speed up the
improvement of SCIP's incumbent (Section~\ref{sec:res-m11}). H2 therefore states that the QUBO pipeline with a
cardinality-preserving sampler outpaces open-source exact solvers, which matters where commercial solvers are not
available; it does not state that the pipeline outpaces the best exact methods. Studies that include
strong commercial solvers find them solving their portfolio problems quickly, but those problems differ from ours (a
binary surrogate, or no cardinality constraint)~\cite{lozano2026where,stopfer2025quantum}; whether H2 would hold against
such solvers is not tested here.

\subsection{Does optimisation accuracy matter out of sample?}\label{sec:disc-oos}

For minimum-variance portfolios, the best non-optimal support is equivalent to the optimum out of sample, and the
out-of-sample variance grows with the in-sample dose along the series of the best alternatives; for mean--variance
portfolios the comparison is inconclusive, and no dose--response is detected at this power
(Section~\ref{sec:res-rq3}). Both statements are about small and specific errors: the best alternatives are a few basis
points of certainty equivalent per year worse than the optimum in sample, and on almost all folds the lower bound of
$K$ holds the weights nearly fixed, so that the alternatives differ mainly in which assets they hold. The difference between the
arms may reflect that the objective of the minimum-variance arm is the variance itself, while the mean--variance
objective also carries the estimated means, whose error is known to dominate out of
sample~\cite{michaud1989markowitz,kan2007optimal,demiguel2009optimal}; with the mean--variance outcome inconclusive,
this is a hypothesis, not a finding. Gilli and Schumann found that beyond some point better solutions of a portfolio
problem do not give better portfolios~\cite{gilli2011optimal,gilli2012heuristic}; constraints and sparsity act as
regularisers~\cite{jagannathan2003risk,brodie2009sparse}. Our design adds certified optima, exactly enumerated
alternatives with measured doses and pre-registered equivalence tests. It confirms equivalence for the minimum-variance
arm and leaves the mean--variance arm open, as the power analysis of the pilot anticipated. On FF49 the real solvers left
almost no in-sample gap, so their own out-of-sample losses could not be calibrated against the curve.

The two parts of the paper meet here. The anytime race of RQ2 on the factor instances is run for start gaps of at
most \res{e2.a3.c2.gstart.max.bp}~bp of certainty equivalent per year, below the smallest effect of interest of RQ3
(25~bp); such in-sample differences between good methods are smaller than the out-of-sample differences that the % num-ok: smallest effect of interest of the plan
mean--variance arm of RQ3 could resolve.

\subsection{Practical guidance}\label{sec:disc-guidance}

For cardinality-constrained mean--variance problems of the sizes studied here, the evidence supports the following
order of work.
\begin{enumerate}
\item Build the common start: solve the continuous relaxation, keep its $K$ largest weights and refine them. It takes a
fraction of a second in the median, and on most of our factor and OR-Library instances it is already within the
target.
\item If a certificate matters, give the start to an exact solver; on the trivial instances the bound closes the
question, at a cost that ranged from seconds at $N = 100$ to more than an hour on some instances
(Section~\ref{sec:res-a3}).
\item To improve on the start, run a local search on the true objective with refinement of the weights; in our
benchmark its median time to the target was within the budget on every instance, and it was the more reliable anytime
method.
\item If a QUBO pipeline is used, for example as the classical counterpart of a quantum sampler, give it a sampler that
preserves the cardinality and a diverse pool: with the penalty samplers of our tuning space it did not improve on
open-source exact solvers (Section~\ref{sec:res-rq2}).
\end{enumerate}

\subsection{Limitations}\label{sec:disc-limits}

The exact baseline is open-source and runs on one thread, its tuning space was small (two post hoc alternatives left
H2 unchanged), and the validation against published times measures certificates and is not normalised for the
hardware; the conclusions of RQ2 are about open-source exact methods. All timings come from one shared
server, with at most 16 pinned single-threaded jobs and the load recorded with each run. The confirmatory cells C1--C3
are synthetic, and their share of trivial instances depends on the design ($K$ half the support of the continuous
optimum, $\lambda = 5$, $\tau = 10^{-3}$); the S\&P cell is descriptive because its pre-registered rule found no
admissible risk aversion, the OR-Library cell has only three port clusters, and the band outcome of H2b depends on it.
References are best known values, not certificates, on \res{e2.n.uncertified} of the \res{e2.n.instances} non-trivial
instances, and most dense instances were left undetermined by the screening. The tuning spaces were frozen before the
pilot: at their penalty levels single-flip annealing could not leave the start, so H1c reads narrowly, and the penalty
samplers were not given smaller penalties; other QUBO samplers, such as parallel tempering or simulated bifurcation,
were not tested. RQ1 is limited to $N \le 24$. The samplers are classical. The ceiling of the surrogate, however, is
the limit that a pipeline with any sampler of the same surrogate approaches as the minimisation improves, quantum
annealers and QAOA included; it is not a bound, since a diverse pool can do better than the exact minimiser of the
surrogate. Out of sample, the primary tests use 49 industry portfolios in a regime of nearly fixed weights; the US
universe is secondary, its calibrated block length was one day, and the power for the mean--variance arm was low.
